\chapter{Conclusion \& Outlook}\label{chap:conclusion}

Resonant soft X-ray scattering has for two decades distinguished the components of organic materials by tuning photon energy across an absorption edge, yet the information it returns about buried molecular states has been bounded by how the optical response is described. This dissertation set out to show that polarized resonant measurements of molecular films already contain more information about those states than prior analyses recover, with the limit lying in the optical model used to interpret them rather than in the measurement. Reflectivity studies of molecular films have obtained that model in one of two ways, tabulating the tensor from NEXAFS of a reference film and holding it fixed at every depth, or fitting it freely at each measured energy. One question was asked of each practice. The first concerned what becomes measurable once the tabulated tensor is replaced by one derived from the electronic transitions of the molecule, and the second concerned how the physical constraints that make such a derivation possible can be carried into the recovery of a tensor that no reference supplies. The approach throughout has been to build the optical tensor from the electronic transitions that produce it, so that each parameter of a fit corresponds to a physical property of the molecule or of its arrangement in the film. Under the first question the tensor was reduced to molecular orientation and density, and under the second the individual resonances themselves were allowed to vary with depth.

For the first question, the ML-refined DFT tensor was built into the analysis of polarized resonant reflectivity from vapor-deposited ZnPc films on native-oxide silicon, with each layer described by a molecular tilt and a density alone. With the spectral shape fixed by the molecule, these few structural parameters must account for the reflectivity at every energy and in both polarizations, and the amplitudes and phases of the fringes across the edge become a joint constraint on them. Across the pre-edge and the resonance series the resulting model described the data as well as or better than free optical constants, from an order of magnitude fewer parameters and with correspondingly smaller uncertainties. Measurements that had been treated as a set of independent energies thereby became a single, overdetermined measurement of molecular state. In a Fisher sensitivity analysis, the constraint on the sub-monolayer interfacial features was traced to the ML-DFT tensor, through the coupling it imposes on fringe amplitude and phase across energy and polarization, and the profile of a second film grown under identical conditions to a different thickness agrees with the first. The first hypothesis therefore holds for ZnPc on silicon, with reflectivity measurements of the kind already in hand fixing the order and density of the interfacial layers, within the assumed depth profile, with uncertainties below the real-space length set by the largest momentum transfer measured.

What an energy-by-energy fit permits became plain in the same comparison. Nothing in such a fit ties the optical constants found at one energy to those found at another. It can return dichroism that no arrangement of the molecule produces, and dispersion inconsistent with the absorption at other energies, a violation of causality that cannot be detected short of sampling the whole edge. In the free fit of ZnPc the surface and buried layers took optical constants that were extreme and poorly determined, including an off-scale dichroism near \qty{287}{eV}, and their low-energy trends broke at the diagnostic energy, leaving the orientation read from them dependent on which energy was analyzed. The ML-DFT tensor is causal and symmetric by construction. Because it admits only ZnPc, the customary assignment of interfacial layers to contamination or oxidation also became testable, since a foreign species at an interface would appear as misfit at the molecule's own resonances. The resonance series of those layers was reproduced with a different tilt and density. Contamination and oxidation are not required by the data within the sensitivity of that series, and the interfacial layers are the molecule itself in other packing and orientational states.

How a wrong tensor behaves in a fit could not be reasoned out in advance, and both of the anticipated behaviors appeared. The largest misfit of the ML-DFT slab model, a phase shift of the s-polarized fringes at the diagnostic energy, was localized to that energy and traced to the structure, since replacing the slabs with the continuous gradient profile removed it with the optical model unchanged. When the same films were fit with a tensor held to the NEXAFS of the reference film, the fit matched the data at the three lower resonant energies and failed at the two energies of strongest orientation sensitivity, placing the misfit where the signature of the error was largest. In Alq\textsubscript{3} a borrowed NEXAFS spectrum misfit every measured energy by a similar factor and was partly absorbed into the structure, with the substrate density driven to a value that the stoichiometry and packing of the molecule cannot allow. We surmise that the difference lies in whether a structural parameter can mimic the error. A change in absorption common to every resonance resembles a density gradient at a fixed spectral shape. Its error spreads across all energies and is partly taken up by the density, whereas an error confined to particular transitions has no structural counterpart and remains as misfit at the energies where those transitions dominate. Only errors of the second kind therefore appear in the residuals of a fixed tensor. Errors of the first kind surface as structural parameters that violate independent constraints, as the Alq\textsubscript{3} density did, and are corrected only when the strengths of individual resonances are free to depart from the reference, the step taken under the second question.

The depth profile returned by this analysis is a product of the first question. At the oxide interface the molecules adopt a lower-density state oriented closer to face-on, and at the free surface a denser state oriented near the magic angle, and each relaxes into the edge-on bulk over roughly one molecular layer. The interfacial state occupies less than a monolayer of a thin film, below the scale that any fit with free optical constants could resolve. Substrate patterning predicts that the first molecular layer already carries the bulk orientation, a prediction inconsistent with a face-on interfacial state beneath an edge-on bulk. The profiles are instead consistent with a growth picture in which substrate adhesion, which favors face-on molecules, competes with π-π cohesion, which favors edge-on packing, so that the gradient records the crossover from interface-dominated to bulk-dominated growth. One consequence is that films grown to partial coverage capture a transient state of the growth front and need not represent the buried interface of a continuously deposited film. These conclusions concern one molecule on one substrate, and the interfacial state on patterning underlayers such as CuI and graphene, where the bulk texture is face-on, can now be measured in the same way.

For the second question, the optical response was written as a sum of Kramers-Kronig consistent oscillator pairs, each a Gaussian absorption with its Dawson-function conjugate tied to an ML-DFT transition cluster and carrying a transition energy, a width, an oscillator strength, and a dipole orientation. Every trial spectrum in a fit is then causal, the tensor follows from the dipoles, and the parameter count is set by the number of transitions in the molecule, where the Kramers-Kronig consistent bases of Kuzmenko and of Stone et al. gain a basis function with every energy measured. With this basis, the amplitudes of individual resonances were refined as functions of depth together with the film structure in both Alq\textsubscript{3} and ZnPc, which satisfies the first part of the second hypothesis. In Alq\textsubscript{3} the degeneracy between the optical model and the film density was removed, the substrate density was no longer driven to an unphysical value, and the misfit fell at every measured energy. The absorption rises toward the substrate across all of the refined resonances, with the third nearly suppressed at the free surface before growing into the bulk and toward the interface. The cause of that evolution cannot yet be assigned. Hydrocarbon contamination at the surface, with the suppressed resonance then an artifact of the model, conformational changes that alter the transition rates near each interface, and hybridization driven by the interfacial potentials all remain consistent with the data, because the optical constants away from the measured energies are inferred from the dispersion of each oscillator.

In the ZnPc films the amplitude of each of the four principal resonances was refined with depth and its dipole cone angle refined once for the whole film, with the tilt and density held at the profile of the first question. The two lowest π* resonances, assigned in the ML-DFT clusters to excitation of two benzo-ring carbon sites into the lowest and a higher unoccupied orbital, exchange oscillator strength between the two interfaces. The first weakens at the free surface and strengthens at the oxide while the second does the reverse, with their summed strength close to the reference value at both boundaries. The third resonance weakens at the oxide interface, the fourth, assigned to the pyrrole carbon bound to nitrogen, strengthens most in relative terms at the oxide where the π* manifold meets the σ* onset, and the cone angles of the π* resonances stay close to the ring normal. These changes differ from one transition to the next, and no rescaling of the reference spectrum reproduces them. The second hypothesis therefore holds for the strengths of individual resonances. We surmise that the exchange between the two lowest resonances is a redistribution of spectral weight between the lowest unoccupied states as the local packing changes. Three mechanisms remain, each with its own signature. Hybridization or charge transfer at the oxide would confine the change to the first molecular layers and shift the transition energies as well as the strengths. Packing-dependent π overlap would make the ratio of the two strengths follow the local tilt wherever the same tilt recurs. Disorder at the free surface would broaden every resonance and lower its anisotropy without exchanging weight between two. Separating them requires a refit with the cone angle of each resonance free at the surface, in the bulk, and at the interface, together with a fixed-angle scan across the edge dense enough to constrain it. Until that measurement is made, the depth dependence of dipole orientation remains open.

The ambiguity in Alq\textsubscript{3} traces to the measurement as much as to the model. Because δ at each energy depends through the Kramers-Kronig relations on β at every other energy, optical constants known exactly at a few energies leave the spectrum between them undetermined, and a handful of energies supplies fewer observables than there are transition parameters to constrain. In synthetic spectra of a known carbon edge, part of the spectrum was recovered with energies placed on the resonances or chosen by a Gaussian-process design, the recovered fraction was essentially unchanged by a tenfold gain in precision at those energies, and the spectrum was constrained far better by a continuous scan across the edge with larger noise at each point than by any sparse set. For sparse measurements the number and placement of energies set the limit before the precision at any one of them does, and the sampling claim of the second hypothesis holds as well. Continuous energy scans should therefore anchor future resonant reflectivity, with the same sampling framework used to choose any additional energies before a beamtime and the oscillator construction serving as a generator of causal optical constants for planning measurements and for training learned inversions.

Two further parts of the second question remain open. The limit in which no calculation supplies the transitions, and only their number is assumed, was argued but not tested on a measured interface. Comparing fits built under different priors against data withheld from the fit could fix the number of oscillators directly, but without the calculation the recovered resonances carry no orbital assignment, and their evolution through a film cannot be traced to its electronic origin. The converse route, refining the ML-DFT clusters themselves against reflectivity with the computed clusters as priors, would correct the misassigned or omitted transitions that a fixed tensor carries into every layer of a film. The energy uncertainties of the present datasets are too large for that refinement, though fixed-angle reflectivity across the edge, measured for the purpose, may suffice.

Every analysis in this dissertation rests on measured NEXAFS, from which the clusters are refined and against which recovered spectra are judged, yet spectra of nominally identical films can differ. An optical model is therefore only as reliable as the spectrum it starts from, and a spectrum must be judged before it is reused. In X-ray Atlas, each raw measurement stays linked to the optical constants derived from it, so that any tensor can be traced to the data and processing that produced it. Every contribution is normalized under one server-side convention, scored by fixed and reproducible metrics of energy sampling, signal to noise, and normalization quality, versioned under a persistent DOI at each change of processing, and credited to each contributor through ORCID. The Python client and command-line tool extend the same records to scripted analysis, and optical models carry the versioning from data into analysis, letting others pin, fork, cite, and reproduce a tensor such as the one in \cref{chap:znpc}. Optical constants thereby become shared and inspectable records of the field, and are no longer files reconstructed separately in each laboratory.

Reflectivity is the one scattering geometry with an exact and inexpensive solution, the one-dimensional 4 × 4 transfer matrix, which made the rigorous optical treatment of \cref{chap:znpc,chap:resonances} tractable. Transmission RSoXS, critical-dimension RSoXS of gratings and lithographic latent images, and grazing-incidence geometries have no equivalent closed form. The same rigor there requires full solutions of Maxwell's equations by finite elements, rigorous coupled-wave analysis, or finite differences in time. A first step in this direction, presented at SPIE Advanced Lithography, recovered the latent chemical profile of a photoresist from critical-dimension RSoXS by fitting finite-element simulations of the full near field in place of a Born approximation, following the approach of Andrle and coworkers.\autocite{andrle2026nanostructure,andrle2025using,li2026pattern}

Neural network surrogates could make these geometries tractable at the scale of a resonant dataset. A direct finite-element solve scales up to \(O\left( N^{3} \right)\) and must be repeated for every energy, polarization, and angle. The open question for the forward problem is whether a trained network, in the manner of MaxwellNet\autocite{lim2022maxwellnet}, can evaluate the scattered field in constant time with accuracy sufficient for fitting. For the inverse problem, the question is whether learned posterior estimators, such as the prior-amortized neural posterior estimation already demonstrated for reflectometry\autocite{starostin2025fast} or joint-embedding predictive architectures\autocite{assran2023self}, can replace iterative fitting with an inversion that costs one network evaluation per measurement.

Closer to the work presented here, beamline 11.0.1.2 is moving to EPICS and Bluesky controls, which will let the decision loop of \cref{fig:3.12} measure until the depth profile, rather than the reflectivity, is determined. The oscillator description of \cref{chap:resonances} extends to polymers such as N2200, whose angle-resolved NEXAFS has already been measured\autocite{gann2014near}, to donor-acceptor blends, and to other absorption edges\autocite{chantler2026nexafs,freychet2022resolving}, and every optical model built this way can be deposited in X-ray Atlas, so that each new analysis begins from a validated, versioned spectrum.

Across these chapters, physical structure was added to the optical model in the form of transition dipoles, Kramers-Kronig consistent oscillators, per-resonance amplitudes, and depth-resolved orientation and density. Each addition replaced freedom that a simpler model had absorbed silently into whichever parameter remained free, and the fit became more tightly constrained even as the description became richer. Buried molecular states become measurable by resonant soft X-ray reflectivity through this structure in the optical model, with no change to the instrument, and the same principle should carry to every geometry in which resonant soft X-rays are used to look inside organic materials.
